\documentclass[11pt]{revtex4-2}

\usepackage[margin=1in]{geometry}
\usepackage{amsmath,amssymb,bm}
\usepackage{physics}
\usepackage{hyperref}



\begin{document}

\title{Full quantum geometric strongly correlated rovibronic dynamics for triatomic molecules: Theory,
 Implementation, and applications}
 \author{Sha Mo}
\author{Bing Gu}
\date{\today}


\begin{abstract}
We present a full quantum  strongly correlated rotation-vibrational-electronic (rovibronic) dynamics for triatomic molecules within the framework of the geometric quantum dynamics. 

The rotational invariance of the potential energy surface and rotational covariance of the electronic states are fully exploited. 

THe critical components are the construction of the nuclear kinetic energy operatro in curvilinear coordiantes,
 and the construction of the electronic overlap matrix in the body-fixed frame using the Eckart conditions. 
 
 In particular, wwe employ the complete active space scf method for the electronic structure calculation and electronic overlap matrix construction.


The implementation is based on the  the discrete variable representation (DVR) and the geometric quantum dynamics. 

A detailed application is presented for the LiH$_2$ molecule to demonstrate the effectiveness of the method. 
 Implementation, and applications to  LiH$_2$
used for full quantum rovibronic wavepacket dynamics.  Starting from the
lab-frame molecular wavefunction, we separate internal coordinates from
overall molecular rotation, transform the electronic basis to a body-fixed
frame, derive the resulting nuclear kinetic energy operator, and obtain the
LDR equation of motion in a fixed total angular momentum basis
\cite{bornhuang1954,bunkerjensen1998,watson1968,zare1988}.  We also state
the approximation made in the current implementation: body-fixed scalar
electronic adiabatic states are used, so explicit electronic angular
momentum gauge terms are neglected.
\end{abstract}
\maketitle

\section{Lab-Frame Molecular Wavefunction}

Let the nuclear configuration of an $N$-atom molecule be described by
Cartesian coordinates $\{\mathbf{R}_A\}$ and masses $\{M_A\}$.  The exact
field-free molecular Coulomb Hamiltonian is
\begin{equation}
  \hat{H}
  =
  \hat{T}_\text{N}
  +
  \hat{H}_{\rm el}(\{\mathbf{R}_A\}),
\end{equation}
with
\begin{equation}
  \hat{T}_{\rm nuc}
  =
  -\frac{1}{2}
  \sum_A \frac{1}{M_A}
  \nabla^2_{A}.
\end{equation}

In the lab frame (LF), the total rovibronic wavefunction can be expanded as a Born-Huang-type ansatz
\cite{bornhuang1954}
\begin{equation}
  \ket{\Psi(t)}
  =
  \sum_a
  \chi_a(\mathbf{q},\Omega,t)
  \ket{\Phi_\alpha(\mathbf{q},\Omega)}
  \label{eq:lab-expansion}
\end{equation}
where $\mathbf{q}$ denotes internal coordinates, $\alpha$ labels the electronic state, and
$
  \Omega = (\alpha,\beta,\gamma)
$
denotes Euler angles specifying the orientation of the molecular body frame
relative to the lab frame.

For a field-free molecule, the electronic Hamiltonian is rotationally
covariant:
\begin{equation}
  \hat{H}_{\rm el}(q,\Omega)
  =
  \hat{U}_{\rm el}(\Omega)
  \hat{H}_{\rm el}(q,0)
  \hat{U}_{\rm el}^\dagger(\Omega).
\end{equation}
Thus the electronic eigenstates at orientation $\Omega$ are rotated copies
of body-fixed electronic states:
\begin{equation}
  \ket{\Phi_a(q,\Omega)}
  =
  \hat{U}_{\rm el}(\Omega)
  \ket{\phi_a(q)} ,
  \label{eq:electronic-rotation}
\end{equation}
where
\begin{equation}
  \hat{H}_{\rm el}(q,0)\ket{\phi_a(q)}
  =
  V_a(q)\ket{\phi_a(q)} .
\end{equation}
where $V_a(q)$ is the potential energy of the electronic state $\alpha$.

The body-fixed electronic representation is obtained by applying the orientation-dependent unitary transformation
\begin{equation}
  \ket{\Psi_{\rm body}}
  =
  \hat{U}_{\rm el}^\dagger(\Omega) \ket{\Psi_{\rm lab}} .
\end{equation}
Using Eq.~\eqref{eq:electronic-rotation},
\begin{equation}
  \ket{\Psi_{\rm body}}
  =
  \sum_a
  \chi_a(q,\Omega,t)
  \ket{\phi_a(q)} .
\end{equation}
The explicit orientation dependence of the electronic basis has therefore
been removed.  It is now carried by the nuclear rotational wavefunction.

The rotational dependence of $\chi_a$ is expanded in Wigner rotational
functions \cite{zare1988,varshalovich1988}:
\begin{equation}
  \chi_a(q,\Omega,t)
  =
  \sum_{J K M}
  C_{J K M,a}(q,t)
  D^J_{M K}(\Omega).
\end{equation}
Equivalently,
\begin{equation}
  \ket{\Psi_{\rm BF}(t)}
  =
  \sum_{J K M a}
  C_{J K M,a}(q,t)
  \ket{q}
  \ket{J K M}
  \ket{\phi_a(q)} .
  \label{eq:body-expansion-continuous}
\end{equation}
Here $J$ is total molecular angular momentum, $M$ is the projection onto the
lab-fixed $Z$ axis, and $K$ is the projection onto the body-fixed $z$ axis.

\section{Coordinate Transformation of the Nuclear Kinetic Energy}

The nuclear Cartesian coordinates can be written as
\begin{equation}
  \mathbf{R}_A
  =
  \mathbf{R}_{\rm cm}
  +
  \mathbf{R}(\Omega)\,
  \mathbf{r}_A(q),
  \label{eq:coordinate-map}
\end{equation}
where $\mathbf{R}_{\rm cm}$ is the center of mass, $\mathbf{R}(\Omega)$ is
the rotation matrix from body to lab frame, and $\mathbf{r}_A(q)$ is the
position of atom $A$ in the body-fixed frame.

After removing translation, the Cartesian Laplacian becomes the
Laplace-Beltrami operator on the curvilinear coordinate space $(q,\Omega)$:
\begin{equation}
  \hat{T}_{\rm nuc}
  =
  -\frac{1}{2}
  |g|^{-1/2}
  \partial_\mu
  \left[
    |g|^{1/2}
    G^{\mu\nu}(q)
    \partial_\nu
  \right],
  \label{eq:laplace-beltrami}
\end{equation}
where $G=g^{-1}$ is the inverse metric and $\mu,\nu = 1, 2, \ldots, 3N-3$ run over internal and
rotational coordinates.

In a body-fixed or Eckart frame, the inverse metric has block structure
\cite{eckart1935,watson1968,bunkerjensen1998,lauvergnat2016}
\begin{equation}
  G
  =
  \begin{pmatrix}
    G_{qq} & G_{q r} \\
    G_{r q} & G_{r r}
  \end{pmatrix},
\end{equation}
where $q$ labels vibrational coordinates and $r$ labels body-fixed
rotations.  Replacing rotational derivatives by body-fixed angular momentum
operators gives the rovibrational kinetic energy
\begin{equation}
  \hat{T}_{\rm rovib}
  =
  \hat{T}_{\rm vib}
  +
  \hat{T}_{\rm rot}
  +
  \hat{T}_{\rm cor}
  +
  V_{\rm extra}(q).
  \label{eq:t-rovib}
\end{equation}

The terms are
\begin{equation}
  \hat{T}_{\rm vib}
  =
  \frac{1}{2}
  \sum_{ij}
  \hat{p}_i
  G_{ij}(q)
  \hat{p}_j,
  \label{eq:t-vib}
\end{equation}
\begin{equation}
  \hat{T}_{\rm rot}
  =
  \frac{1}{2}
  \sum_{\alpha\beta}
  G_{\alpha\beta}(q)
  \hat{J}_\alpha
  \hat{J}_\beta,
  \label{eq:t-rot}
\end{equation}
and
\begin{equation}
  \hat{T}_{\rm cor}
  =
  \frac{1}{2}
  \sum_{i\alpha}
  \left[
    \hat{p}_i G_{i\alpha}(q)
    +
    G_{\alpha i}(q) \hat{p}_i
  \right]
  \hat{J}_\alpha.
  \label{eq:t-cor}
\end{equation}
Here $\hat{p}_i=-i\partial/\partial q_i$, and
$\hat{J}_\alpha$ are body-fixed angular momentum operators.  The scalar
$V_{\rm extra}(q)$ is the curvilinear metric or pseudopotential term caused
by the volume element and operator ordering \cite{watson1968,lauvergnat2016}.

\section{Effect of the Electronic Body-Frame Unitary on the KEO}

The transformation
\begin{equation}
  \hat{H}_{\rm body}
  =
  \hat{U}_{\rm el}^\dagger(\Omega)
  \hat{H}_{\rm lab}
  \hat{U}_{\rm el}(\Omega)
\end{equation}
does not leave the nuclear kinetic energy unchanged, because
$\hat{T}_{\rm nuc}$ contains derivatives with respect to nuclear coordinates,
including derivatives with respect to $\Omega$:
\begin{equation}
  \hat{U}_{\rm el}^\dagger
  \partial_\Omega
  \hat{U}_{\rm el}
  \neq
  \partial_\Omega .
\end{equation}
Instead, derivatives become covariant derivatives:
\begin{equation}
  \partial_\Omega
  \rightarrow
  \partial_\Omega
  +
  \hat{U}_{\rm el}^\dagger
  (\partial_\Omega \hat{U}_{\rm el}) .
\end{equation}
The connection is related to electronic angular momentum:
\begin{equation}
  \hat{U}_{\rm el}^\dagger
  (\partial_\Omega \hat{U}_{\rm el})
  \sim
  -i \hat{\mathbf{j}}_{\rm el}.
\end{equation}
Therefore, in the most general body-fixed theory, the rotational angular
momentum entering the KEO is shifted schematically as
\begin{equation}
  \hat{\mathbf{J}}
  \rightarrow
  \hat{\mathbf{J}} - \hat{\mathbf{j}}_{\rm el}.
  \label{eq:j-shift}
\end{equation}
For example,
\begin{equation}
  \hat{T}_{\rm rot}
  \sim
  \frac{1}{2}
  \sum_{\alpha\beta}
  G_{\alpha\beta}(q)
  \left(
    \hat{J}_\alpha - \hat{j}^{\rm el}_\alpha
  \right)
  \left(
    \hat{J}_\beta - \hat{j}^{\rm el}_\beta
  \right).
\end{equation}

The present implementation assumes scalar, body-fixed adiabatic electronic
states with quenched explicit electronic angular momentum, so
\begin{equation}
  \hat{\mathbf{j}}_{\rm el} \approx 0.
\end{equation}
Under this approximation, Eqs.~\eqref{eq:t-vib}--\eqref{eq:t-cor} are used
with $\hat{\mathbf{J}}$ only.

\section{Local Diabatic Representation on a DVR Grid}

Let the internal coordinates be represented by DVR grid points $q_n$.  The
body-fixed LDR expansion is built on the DVR idea that coordinate-local
grid functions provide an efficient representation of nuclear kinetic
operators and local potentials \cite{lightcarrington2000}.  The expansion is
\begin{equation}
  \ket{\Psi(t)}
  =
  \sum_{n \tau a}
  C_{n\tau a}(t)
  \ket{q_n}
  \ket{J\tau}
  \ket{\phi_a(q_n)} ,
  \label{eq:ldr-expansion}
\end{equation}
where
\begin{equation}
  \tau = (K,M)
\end{equation}
for the full $|J K M\rangle$ basis, while
\begin{equation}
  \tau = K
\end{equation}
for a fixed-$M$ block.

The local electronic overlap matrix is
\begin{equation}
  A_{n a, m b}
  =
  \braket{\phi_a(q_n)}{\phi_b(q_m)} .
  \label{eq:electronic-overlap}
\end{equation}
This overlap must be evaluated in a consistent body-fixed or Eckart gauge.
If electronic states are computed in inconsistent lab-frame orientations,
$A$ contains spurious rotational gauge dependence.

The electronic potential energy is local and diagonal:
\begin{equation}
  \bra{\phi_a(q_n)}
  \hat{H}_{\rm el}(q_n)
  \ket{\phi_b(q_n)}
  =
  E_a(q_n)\delta_{ab}.
\end{equation}

The rovibrational KEO matrix is
\begin{equation}
  T^J_{n\tau,m\tau'}
  =
  \mel{q_n,J\tau}
  {\hat{T}_{\rm rovib}}
  {q_m,J\tau'} .
  \label{eq:tj-matrix}
\end{equation}

Projecting the time-dependent Schrodinger equation gives
\begin{equation}
  i\dot{C}_{n\tau a}
  =
  E_a(q_n) C_{n\tau a}
  +
  \sum_{m\tau' b}
  T^J_{n\tau,m\tau'}
  A_{n a,m b}
  C_{m\tau' b}.
  \label{eq:ldr-eom-compact}
\end{equation}
This is the compact LDR rovibronic equation of motion.

If one separates the derivative part of the KEO from the scalar
pseudopotential, Eq.~\eqref{eq:ldr-eom-compact} may be written more
explicitly as
\begin{align}
  i\dot{C}_{n\tau a}
  &=
  \left[
    E_a(q_n) + V_{\rm extra}(q_n)
  \right]
  C_{n\tau a}
  \nonumber\\
  &\quad+
  \sum_{m\tau' b}
  T^{J,{\rm deriv+rot+cor}}_{n\tau,m\tau'}
  A_{n a,m b}
  C_{m\tau' b}.
  \label{eq:ldr-eom-explicit}
\end{align}

\section{Special Cases}

\subsection{$J=0$}

For $J=0$, there is only one rotational state:
\begin{equation}
  K = M = 0,
  \qquad
  n_{\rm rot}=1.
\end{equation}
All angular momentum matrices vanish:
\begin{equation}
  \hat{J}_x = \hat{J}_y = \hat{J}_z = 0.
\end{equation}
Therefore,
\begin{equation}
  \hat{T}_{\rm rot}=0,
  \qquad
  \hat{T}_{\rm cor}=0,
\end{equation}
and the dynamics reduces to vibronic dynamics with the curvilinear
vibrational KEO:
\begin{equation}
  \hat{T}_{\rm nuc}
  =
  \hat{T}_{\rm vib}
  +
  V_{\rm extra}.
\end{equation}

\subsection{Fixed $J_z=M$}

For field-free dynamics, the lab-frame projection $M$ is conserved.  One may
therefore work in a fixed-$M$ block:
\begin{equation}
  \ket{J K M_0},
  \qquad
  K=-J,\ldots,J.
\end{equation}
The rotational dimension is then
\begin{equation}
  n_{\rm rot} = 2J+1,
\end{equation}
instead of the full
\begin{equation}
  n_{\rm rot} = (2J+1)^2.
\end{equation}
This reduction is valid for field-free rovibronic dynamics.  In the presence
of external fields, whether $M$ remains conserved depends on the polarization
and selection rules of the light-matter interaction.

\section{Implementation details}

The current PyQED implementation uses the following assumptions:
\begin{enumerate}
  \item Electronic states are scalar body-fixed adiabatic states
  $\ket{\phi_a(q)}$.
  \item Electronic overlaps $A_{na,mb}$ are computed in a consistent
  body-fixed or Eckart orientation.
  \item Explicit electronic angular momentum gauge terms
  $\hat{\mathbf{j}}_{\rm el}$ are neglected.
  \item The full rovibrational KEO includes vibrational, rotational,
  Coriolis, and metric pseudopotential terms.
  \item Setting $J=0$ gives vibronic dynamics.  Setting $J>0$ gives
  rovibronic dynamics.  Supplying a fixed $J_z=M$ exploits conservation of
  lab-frame angular momentum projection in field-free dynamics.
\end{enumerate}










\end{document}
